%!TEX program = lualatex
% SPDX-License-Identifier: CC-BY-SA-4.0
% Requires: https://github.com/jemmybutton/fiziko (GPL-3.0)
% Set MPINPUTS to the directory containing fiziko.mp before running lualatex.
\documentclass[11pt]{article}

\usepackage{iftex}
\ifLuaTeX\else
  \PackageError{vintage-latex}{This example requires LuaLaTeX}
    {Run lualatex, not pdflatex or xelatex.}
\fi

\usepackage{fontspec}
\usepackage{unicode-math}
\usepackage{xcolor}
\usepackage{geometry}
\usepackage[tracking=true,protrusion=false,expansion=false]{microtype}
\usepackage{luamplib}

\setmainfont{EB Garamond}[Numbers={OldStyle,Proportional}]
\setmathfont{Garamond-Math.otf}
\definecolor{ragpaper}{HTML}{F3E9D6}
\definecolor{printerink}{HTML}{25201B}
\pagecolor{ragpaper}
\color{printerink}
\geometry{a4paper,landscape,top=18mm,bottom=18mm,left=24mm,right=24mm}
\pagestyle{empty}
\setlength{\parindent}{0pt}

\mplibforcehmode
\everymplib{
  input fiziko.mp;
  randomseed := 1729;
  defineMinStrokeWidth(0.22pt);
  defineLightDirection(-1/8pi, 1/8pi);
  color inkcolor; inkcolor := (0.145, 0.125, 0.106);
  drawoptions(withcolor inkcolor);
  beginfig(1);
}
\everyendmplib{endfig;}

\begin{document}

\begin{center}
  {\small\textls[120]{\textsc{Plate II}}}\par
  \vspace{.3em}
  {\Large The Astronomical Telescope of Kepler}\par
  \vspace{.2em}
  {\itshape the course of the rays traced through two convex glasses}
\end{center}

\vspace{1.2em}

\begin{center}
\begin{mplibcode}
  numeric nGlass, eyepieceAt;
  path objective, eyepiece, axis, ray[];
  pair focus, exitPupil, a_, b_;

  nGlass := 1.52;
  eyepieceAt := 15.6cm;
  % Radii of curvature in cm, positive for convex; thickness; diameter.
  objective := lens((14, 14), 0.5, 3.6, cm);
  eyepiece := lens((2.2, 2.2), 0.32, 1.4, cm) shifted (eyepieceAt, 0);
  axis := (-3.4cm, 0) -- (eyepieceAt + 4.6cm, 0);

  draw axis dashed evenly scaled 0.6 withpen pencircle scaled 0.28pt;
  draw solid(objective, 52, 1);
  draw objective withpen pencircle scaled 0.6pt;
  draw solid(eyepiece, 52, 1);
  draw eyepiece withpen pencircle scaled 0.6pt;

  % The exit ray of refractionPath is short; extend it along its direction.
  vardef extended(expr p, l) =
    p -- (point length(p) of p
          + l * unitvector(point length(p) of p - point (length(p)-1) of p))
  enddef;

  % Bundle from a star on the axis: parallel rays, brought to the focus F,
  % then made parallel again by the eyepiece.
  for h = -1.2cm, -0.6cm, 0.6cm, 1.2cm:
    ray1 := (-3.2cm, h) -- (-1.6cm, h);
    ray2 := refractionPath(ray1, objective, nGlass);
    ray3 := refractionPath(ray2, eyepiece, nGlass);
    draw extended(ray3, 3.4cm) withpen pencircle scaled 0.32pt;
  endfor;
  ray1 := (-3.2cm, 1.2cm) -- (-1.6cm, 1.2cm);
  ray2 := refractionPath(ray1, objective, nGlass);
  a_ := point (length(ray2)-1) of ray2;  b_ := point length(ray2) of ray2;
  focus = whatever[a_, b_];  ypart focus = 0;

  % Bundle from a second star two degrees off the axis, drawn broken.
  for h = -1.2cm, 0, 1.2cm:
    ray1 := ((-3.2cm, h) -- (-1.6cm, h)) shifted (0, 3.2cm*sind(2)) rotated -2;
    ray2 := refractionPath(ray1, objective, nGlass);
    ray3 := refractionPath(ray2, eyepiece, nGlass);
    draw extended(ray3, 3.4cm) dashed evenly scaled 0.5
      withpen pencircle scaled 0.32pt;
  endfor;
  % Its chief ray crosses the axis at the exit pupil, where the eye is placed.
  ray1 := ((-3.2cm, 0) -- (-1.6cm, 0)) shifted (0, 3.2cm*sind(2)) rotated -2;
  ray2 := refractionPath(ray1, objective, nGlass);
  ray3 := refractionPath(ray2, eyepiece, nGlass);
  a_ := point (length(ray3)-1) of ray3;  b_ := point length(ray3) of ray3;
  exitPupil = whatever[a_, b_];  ypart exitPupil = 0;

  eyescale := 1.1cm;
  draw eye(180) shifted exitPupil shifted (1.05cm, 0);

  % Lettering, as on the engraved plates: capitals near the parts.
  label.top(btex $A$ etex, (0.25cm, 1.95cm));
  label.top(btex $B$ etex, (eyepieceAt + 0.16cm, 0.85cm));
  dotlabel.bot(btex $F$ etex, focus);
  dotlabel.bot(btex $O$ etex, exitPupil);
  label.bot(btex \small\itshape rays from a star upon the axis etex,
    (-1.9cm, -1.55cm));
  label.bot(btex \small\itshape rays from a star two degrees from it etex,
    (-1.9cm, -2.05cm));
  draw (-3.2cm, -1.55cm) -- (-4.3cm, -1.55cm) withpen pencircle scaled 0.32pt;
  draw (-3.2cm, -2.05cm) -- (-4.3cm, -2.05cm) dashed evenly scaled 0.5
    withpen pencircle scaled 0.32pt;
\end{mplibcode}
\end{center}

\vspace{1.2em}
\begin{minipage}{.48\linewidth}
\small
The object-glass $A$ receives the rays of a star nearly parallel, and unites
them at its focus $F$. Rays from a second star, inclined to the first, are
brought to a focus a little above or below $F$, so that in the focal plane
there lies a small inverted image of the field.
\end{minipage}\hfill
\begin{minipage}{.48\linewidth}
\small
The eye-glass $B$ is set at its own focal distance beyond that plane. It
returns each pencil to parallelism, but the pencils now diverge from one
another more steeply than the stars themselves; the whole of them pass
through the small opening at $O$, where the pupil of the observer is placed.
\end{minipage}

\vspace{1em}
\footnotesize
No ray in this plate is drawn to a predetermined point. Each is refracted at
every glass surface by fiziko's tracing macro, the focus and the exit pupil
are found where computed rays cross the axis, and the lettering shares the
fonts of the page. Changing a radius of curvature moves every ray, the letter
$F$, and the eye together.

\end{document}
